%! Singular Values and Singular Vectors — Unit 7, Lesson 7.1 — Group Activity
\documentclass[10pt]{article}
\usepackage{linalg-article}
\usepackage{linalg-boxes}
\newcommand{\TallMath}[1]{$\displaystyle #1\rule[-1.4em]{0pt}{3.2em}$}
\newcommand{\uu}{\mathbf{u}}
\newcommand{\vv}{\mathbf{v}}
\newcommand{\ee}{\mathbf{e}}
\newcommand{\zero}{\mathbf{0}}
\newcommand{\T}{^{\top}}

\begin{document}

\pageheader{Unit 7, Lesson 7.1}{Group Activity}
\namepartnerperiod

\vspace{0.10in}

\begin{headlinebox}{blush}
\small\textbf{Finding the SVD with $A\T A$.} The recipe: \textbf{(1)} build $A\T A$; \textbf{(2)} find its
eigenvalues $\lambda_1\ge\lambda_2\ge0$ and perpendicular unit eigenvectors $\vv_i$; \textbf{(3)}
$\sigma_i=\sqrt{\lambda_i}$; \textbf{(4)} $\uu_i=A\vv_i/\sigma_i$; \textbf{(5)} assemble $A=U\Sigma V\T$.
Work with your partner, show every step, and \emph{justify} each answer.
\end{headlinebox}

\vspace{0.10in}

\begin{tcolorbox}[colback=white, colframe=black!40, title=\textbf{Tier R --- Remediate},
    fonttitle=\bfseries, arc=1.5mm, left=3mm, right=3mm, top=2mm, bottom=2mm, breakable]
When the columns of $A$ are perpendicular, $A\T A$ is \emph{diagonal} --- so you can read $\lambda$ right
off, and the input vectors are the axes $\ee_1,\ee_2$.
\begin{enumerate}[leftmargin=1.7em, itemsep=9pt]
  \item \textbf{Diagonal.} For $A=\begin{bsmallmatrix} 3 & 0\\ 0 & 2 \end{bsmallmatrix}$, compute
        $A\T A=\blank{2.2cm}$. Its eigenvalues are $\lambda_1=\blank{0.9cm}$, $\lambda_2=\blank{0.9cm}$,
        so $\sigma_1=\blank{0.8cm}$, $\sigma_2=\blank{0.8cm}$. With $\vv_1=\ee_1$, $\vv_2=\ee_2$, find
        $\uu_1=A\vv_1/\sigma_1$ and $\uu_2=A\vv_2/\sigma_2$. \writeline
  \item \textbf{Swap --- watch the order.} For $S=\begin{bsmallmatrix} 0 & 3\\ 2 & 0 \end{bsmallmatrix}$,
        compute $S\T S$ (still diagonal). Its eigenvalues are $9$ and $4$. \emph{Which axis} is the
        eigenvector for $\lambda=9$? Give $\sigma_1,\sigma_2$ and then $\uu_1=S\vv_1/\sigma_1$. Is $\uu_1$
        the same direction as $\vv_1$? \writelines{2}
\end{enumerate}
\end{tcolorbox}

\begin{tcolorbox}[colback=white, colframe=black!40, title=\textbf{Tier A --- Approaching Proficiency},
    fonttitle=\bfseries, arc=1.5mm, left=3mm, right=3mm, top=2mm, bottom=2mm, breakable]
Now a genuinely tilted matrix --- the columns are \emph{not} perpendicular, so $A\T A$ is not diagonal.
Run the full recipe on $A=\begin{bsmallmatrix} 2 & 2\\ 1 & -2 \end{bsmallmatrix}$.
\begin{enumerate}[leftmargin=1.7em, itemsep=9pt]
  \item \textbf{Steps 1--3.} Compute $A\T A$, find its eigenvalues $\lambda_1\ge\lambda_2$, and take
        $\sigma_1,\sigma_2$. \writelines{2}
  \item \textbf{Input vectors.} The eigenvector directions are $\begin{bsmallmatrix} 1\\ 2 \end{bsmallmatrix}$
        (for $\lambda_1$) and $\begin{bsmallmatrix} 2\\ -1 \end{bsmallmatrix}$ (for $\lambda_2$). Normalize
        each to a unit $\vv_i$ (divide by $\sqrt5$), and confirm $\vv_1\perp\vv_2$. \writeline
  \item \textbf{Output vectors + check.} Compute $\uu_1=A\vv_1/\sigma_1$ and $\uu_2=A\vv_2/\sigma_2$, and
        verify $\uu_1\cdot\uu_2=0$ --- perpendicular out, automatically. \writelines{2}
\end{enumerate}
\end{tcolorbox}

\begin{tcolorbox}[colback=white, colframe=black!40, title=\textbf{Tier E --- Extension},
    fonttitle=\bfseries, arc=1.5mm, left=3mm, right=3mm, top=2mm, bottom=2mm, breakable]
\textbf{Any shape.} The recipe needs only $A\T A$, so $A$ need not be square. Take the tall
$3\times2$ matrix $A=\begin{bsmallmatrix} 2 & 0\\ 0 & 2\\ 1 & 2 \end{bsmallmatrix}$ --- it has no
eigenvalues of its own (it isn't square), but $A\T A$ is a clean $2\times2$.
\begin{enumerate}[leftmargin=1.7em, itemsep=9pt]
  \item Compute $A\T A$ (a $2\times2$: each entry is a dot product of two columns of $A$). \writelines{2}
  \item Find its eigenvalues and take $\sigma_1,\sigma_2$. \writeline
  \item \textbf{Justify.} Explain why the eigenvalues of \emph{any} $A\T A$ must be $\ge0$ --- so that
        $\sigma=\sqrt{\lambda}$ always exists. (Hint: what does $\lambda$ equal?) \writeline
\end{enumerate}
\end{tcolorbox}

\end{document}
